% Part: many-valued-logic
% Chapter: infinite-valued-logics
% Section: introduction

\documentclass[../../../include/open-logic-section]{subfiles}

\begin{document}

\olfileid[id]{mvl}{inf}{int}

\olsection{Pengantar}

Banyaknya nilai kebenaran dalam suatu matriks tidak harus berhingga. Pilihan
yang jelas untuk himpunan nilai kebenaran yang tak berhingga ialah himpunan
bilangan rasional antara $0$ dan~$1$, $V_\infty = [0,1] \cap \Rat$, yakni,
\begin{align*}
    V_\infty & = \Setabs{\frac{n}{m}}{n,m \in \Nat \text{ dan } n\le m}.
\intertext{Ketika mempertimbangkan himpunan nilai kebenaran tak hingga ini,
sering kali berguna untuk juga mempertimbangkan himpunan-himpunan bagiannya}
V_m & = \Setabs{\frac{n}{m-1}}{n \in \Nat \text{ dan } n < m}
\intertext{Sebagai contoh, $V_5$ merupakan himpunan dengan $5$ nilai kebenaran
yang berjarak sama,}
V_5 & = \{0, \frac{1}{4}, \frac{1}{2}, \frac{3}{4}, 1\}.
\end{align*}
Dalam logika yang didasarkan pada himpunan-himpunan nilai kebenaran ini,
biasanya $1$ merupakan satu-satunya nilai kebenaran yang ditetapkan, yakni
$V^+ = \{1\}$. Dengan kata lain, kita membiarkan $1$ memainkan peran kebenaran
(mutlak) dan $0$ memainkan peran kepalsuan mutlak, tetapi !!{formula}s dapat
mengambil nilai perantara mana pun dalam~$V$.

Namun, kita juga dapat mempertimbangkan himpunan $V_{[0,1]} = [0,1]$ dari
semua bilangan \emph{real} antara $0$ dan~$1$, atau himpunan bagian tak
berhingga lainnya dari~$[0,1]$. Logika dengan himpunan nilai kebenaran ini
sering disebut \emph{logika kabur}.


\end{document}
